\chapter{前沿假说与证据边界}
\label{app:frontier}

\section{本附录的目的}

本附录的存在，是为了保护本书的可信度。有些关于 consciousness 的提案
富有想象力，也可能具有重要性，但它们目前还不属于与经验神经科学或
成熟动力系统工具相同的证据层级。本书并不忽略它们，而是将其隔离出来。

\section{为什么隔离很重要}

\begin{enumerate}[nosep]
  \item 它能防止经验核心因为思辨性越界而被削弱。
  \item 它能让有雄心的读者看到前沿思考正在朝哪里推进。
  \item 如果其中一些理论后来走向成熟，它能为本书后续版本提供一个
        清晰的扩展位置。
\end{enumerate}

\section{当前前沿示例}

一个前沿路线的例子，是围绕意识架构展开的 resonance-complexity 式理论化。
这类工作也许对未来整合有用，但在当前这本书里，它应当被视为探索性的：
\begin{itemize}[nosep]
  \item 它可以刺激概念比较；
  \item 它可以推动新的数学问题；
  \item 它不应被用来替代强有力的经验性证据。
\end{itemize}

\section{阅读指引}

\begin{readingbox}
\textbf{基础}
\begin{itemize}[nosep]
  \item Christof Koch, \emph{Consciousness}.
  \item Jesse Prinz, \emph{The Conscious Brain}.
\end{itemize}
\textbf{现代研究}
\begin{itemize}[nosep]
  \item 在接触前沿提案之前，先阅读经验性意识章节。
\end{itemize}
\textbf{前沿}
\begin{itemize}[nosep]
  \item \emph{Resonance Complexity Theory and the Architecture of
        Consciousness} (arXiv, 2025).
\end{itemize}
\end{readingbox}

\begin{summarybox}
本附录是一个治理机制。它让本书能够承认前沿想法的存在，
同时又不会悄悄把它们提升为主要证据骨架的一部分。
\end{summarybox}